%% Jarkom report using its original Word-matched layout
\documentclass[12pt]{article}
\usepackage{config}

\begin{document}
\pagenumbering{arabic}\setcounter{page}{0}

%% ==================================================================
%%  HALAMAN SAMPUL (halaman 0)
%% ==================================================================
\thispagestyle{empty}
{\linespread{1}\selectfont\begin{tikzpicture}[remember picture,overlay,
      shift={(current page.north west)},x=0.354mm,y=-0.354mm]
    % teks kiri atas (Word: Cambria 14pt)
    \node[anchor=north west,inner sep=0] at (49.2,45.2) {\cambriafont\fontsize{14}{16}\selectfont Laporan Praktikum};
    % kotak "Pertemuan" (Word PartTitle Arial Black 10pt, angka PartLabel 52pt)
    \fill[abuabu] (378.5,47.8) rectangle (448.6,128.3);
    \node[anchor=north,inner sep=0] at (413.5,51)
    {\arialblack\fontsize{10}{11}\selectfont Pertemuan};
    \node[anchor=north,inner sep=0,text=white] at (413.5,82)
    {\arialblack\fontsize{52}{52}\selectfont\pertemuan};
    % kode mata kuliah + garis (Word CompanyName: Garamond)
    \node[anchor=north west,inner sep=0] at (48.6,148.6)
    {\fontsize{12.3}{14}\selectfont IF25-21008: JARINGAN KOMPUTER};
    \draw[line width=0.5pt] (45.4,162.5) -- (449.4,162.5);
    \node[anchor=north west,inner sep=0] at (48.6,170)
    {\fontsize{18.5}{21}\selectfont Laporan Praktikum/Lab};
    % banner judul (Word TitleCover: Bernard MT Condensed 24pt Bold)
    \fill[abubanner] (15.1,236.7) rectangle (479.7,333.9);
    \node[anchor=center,inner sep=0,align=center,font=\bernard\bfseries\fontsize{24}{28}\selectfont] at (249,287.6)
    {\judul};
    % identitas (rata kanan, Word: Garamond 14pt Bold)
    \node[anchor=north east,inner sep=0,align=right,font=\bfseries\fontsize{14}{15.6}\selectfont] at (449.2,379.2)
    {\namapeserta\\ \kelas\\ \nim};
    \node[anchor=north east,inner sep=0,align=right,font=\bfseries\fontsize{14}{15.6}\selectfont] at (449.2,444.4)
    {\underline{Tim Asisten}\\ \asistenA\\ \asistenB};
    % institusi (Word: Garamond 14pt Bold)
    \node[anchor=north east,inner sep=0,align=right,font=\bfseries\fontsize{14}{15.4}\selectfont] at (449.2,545.8)
    {Teknik Informatika\\ Fakultas Teknologi Industri\\
      Institut Teknologi Sumatera\\ \bulantahun};
    % logo institusi (tiga logo, rata bawah)
    \node[anchor=base west,inner sep=0] at (42,649.4)
    {\raisebox{-2.4pt}{\includegraphics[width=57.6pt]{figure/logo-itera.png}}\hspace{23pt}%
      \includegraphics[width=114.3pt]{figure/logo-fti.png}\hspace{20.3pt}%
      \includegraphics[width=156.9pt]{figure/logo-if.png}};
\end{tikzpicture}}
\newpage

%% ==================================================================
%%  HALAMAN JUDUL DALAM (halaman 1)
%% ==================================================================
\thispagestyle{halamandua}
{\linespread{1}\selectfont\begin{tikzpicture}[remember picture,overlay,
      shift={(current page.north west)},x=0.354mm,y=-0.354mm]
    \node[anchor=north,inner sep=0,align=center,text=abujudul,font=\fontsize{9.4}{16.6}\selectfont\addfontfeatures{LetterSpace=14.5}] at (248.6,139.6)
    {FAKULTAS TEKNOLOGI INDUSTRI\\ ITERA 2026};
    % judul boleh dipecah dengan \\; garis bawah ikut turun mengikuti baris terakhir
    \node[name=jd,anchor=north,inner sep=0,align=center,text=abujudul,font=\arialblack\fontsize{23.5}{27}\selectfont\addfontfeatures{LetterSpace=-7}] at (251.8,206.2)
    {\judul\rule[-8pt]{0pt}{8pt}};
    \draw[abujudul,line width=0.6pt] (70,0 |- jd.south) ++(0,12.44) -- ++(356.4,0);
    % anggota kelompok (Word ReturnAddress: Garamond Bold), tumbuh ke atas
    % kosongkan \daftaranggota untuk laporan individu agar blok ini tidak tampil
    \ifx\daftaranggota\empty\else
      \node[anchor=south,inner sep=0,align=center,font=\bfseries\fontsize{9.6}{11.2}\selectfont] at (249,527)
      {Nama anggota kelompok:\\\daftaranggota
        \ifx\catatananggota\empty\else\\[6pt]\catatananggota\fi};
    \fi
    % alamat institusi (Word ReturnAddress: Garamond)
    \node[anchor=north,inner sep=0,align=center,font=\fontsize{9.6}{11.2}\selectfont] at (249,541.3)
    {ITERA Laboratory\\ Program Studi Teknik Informatika\\
      Fakultas Teknologi Industri\\ Institut Teknologi Sumatera};
    \node[anchor=north,inner sep=0,align=center,font=\fontsize{9.6}{11.2}\selectfont] at (249,596.7)
    {Jl. Terusan Ryacudu, Desa Way Hui, Kecamatan Jati Agung,\\
      Lampung Selatan -- Lampung 35365\\
      Telp. +62-721-8030188; 8030189};
\end{tikzpicture}}
\newpage
%% ==================================================================
%%  ISI LAPORAN (halaman 2 dst.)
%% ==================================================================
\pagestyle{isi}

% Semua bagian laporan digabung di satu berkas.
% Ganti setiap \titik dengan poin rangkuman praktikum.
% Word Heading2b: Arial bold italic, isi Garamond dengan bullet kotak.
\vspace*{-8.8mm}
{\centering\bfseries\fontsize{15}{18}\selectfont\judul\par}\vspace{1.4mm}
\begin{rangkuman}
  \linespread{1.0}\selectfont
  {\arialfont\bfseries\itshape\fontsize{12}{14}\selectfont Rangkuman praktikum}\par\vspace{4mm}
  {\arialfont\bfseries\itshape\fontsize{10}{12}\selectfont
  Poin-poin utama yang anda ketahui selama praktikum adalah:\par}\vspace{-1.5mm}
  \begin{itemize}[label=\raisebox{0.45ex}{\fontsize{3.2}{3.2}\selectfont$\blacksquare$},leftmargin=10.9mm,labelsep=3.85mm,labelwidth=0pt,
                  align=left,itemsep=0pt,topsep=0pt,parsep=0pt,
                  before=\normalfont\fontsize{9.6}{17.4}\selectfont]
    \item \titik
    \item \titik
    \item \titik
    \item \titik
    \item \titik
  \end{itemize}
\end{rangkuman}
\section{Dokumentasi Praktikum}

\subsection{Jelaskan Langkah-langkah yang dilakukan selama praktikum}

\begin{itemize}
  \item Langkah pertama adalah (silahkan jelaskan langkah-langkah
      selama praktikum)

      Lampirkan gambar yang dilakukan selama langkah pertama.

      \begin{figure}[H]
        \centering
        \gambarpengganti{figure/contoh-gambar.png}
        \caption{Ilustrasi dokumentasi praktikum}
      \end{figure}
  \newpage % sama dengan Word: langkah kedua mulai di halaman baru (hapus jika tidak perlu)
  \item Langkah kedua adalah (silahkan jelaskan langkah-langkah
      selama praktikum)

      Jika ada tabel yang perlu dilampirkan berikut format table nya.

      \begin{table}[H]
        \centering
        \caption{Contoh penggunaan table}
        \linespread{1}\fontsize{8}{10}\selectfont
        \rowcolors{2}{tabband}{white}
        \setlength{\tabcolsep}{3pt}
        \renewcommand{\arraystretch}{0.9}
        \arrayrulecolor{tabline}
        \begin{tabular}{!{\color{tabline}\vrule}>{\centering\arraybackslash}p{9mm}!{\color{tabline}\vrule}p{30mm}!{\color{tabline}\vrule}p{30mm}!{\color{tabline}\vrule}p{30mm}!{\color{tabline}\vrule}}
          \hline
          \rowcolor{tabhead}
          \thead{No} & \thead{Judul kolom 1} & \thead{Judul Kolom 2} & \thead{Judul Kolom 3} \\
          \hline
          & & & \\ \hline
          & & & \\ \hline
          & & & \\ \hline
          & & & \\ \hline
          & & & \\ \hline
          & & & \\ \hline
          & & & \\ \hline
          & & & \\ \hline
        \end{tabular}
      \end{table}
  \item Langkah ketiga adalah (silahkan jelaskan langkah-langkah
      selama praktikum). Jika ada penggunaan list atau item di dalam
      penjelasan anda, silahkan gunakan format bullet and numbering berikut:
      \begin{itemize}
        \item List 1
        \item List 2
        \item List 3
      \end{itemize}
\end{itemize}
\par
\subsubsection{Alat dan Bahan}

Jelaskan alat dan bahan yang digunakan dalam praktikum. Deskripsikan alat atau
device yang anda gunakan. Jika menggunakan komputer Windows, cukup dengan
memberikan screenshoot dxdiag.
\par
\par
\subsubsection{Tugas}

Silahkan jelaskan jawaban dari tugas yang diberikan. Jika tidak ada tugas, maka
setiap praktikum anda wajib melampirkan dokumen laporan praktikum ini.
\par
\nocite{forouzan2007}






\newpage
\judulreferensi{Referensi}
\nocite{*}
\printbibliography[heading=none]

\end{document}
